\documentclass[14pt]{extreport}

\usepackage[margin=1in]{geometry}
\usepackage[T1]{fontenc}
\usepackage[utf8]{inputenc}
\usepackage{hyperref}

\begin{document}

\begin{center}
\textbf{Experiment 7: MENACE and Multi-Armed Bandits}\\[0.8em]
Manoj Kumar Panda -- 20251602013\\
Meetkumar Kankotiya -- 20251602007\\
Shrey Parikh -- 20251602017\\[0.5em]
Group--10\\
CS663 -- Artificial Intelligence\\
Indian Institute of Information Technology Vadodara (IIITV)
\end{center}

\vspace{1em}

\noindent\textbf{Abstract:}\\
This experiment studies reinforcement learning through two core setups:
multi-armed bandits and the Matchbox Educable Naughts and Crosses Engine
(MENACE). Bandit problems are used to illustrate the exploration--exploitation
tradeoff and non-stationary rewards, while MENACE demonstrates how a simple
probabilistic mechanism can learn to play Tic-Tac-Toe. The Python
implementation uses $\epsilon$-greedy agents for bandits and a bead-based
policy for MENACE.

\vspace{1em}

\noindent\textbf{Introduction:}\\
Reinforcement learning (RL) is concerned with how an agent learns to act by
interacting with an environment and receiving feedback in the form of rewards.
Two classic RL examples are multi-armed bandits and the MENACE system. Bandits
capture the core exploration--exploitation dilemma in a very simple setting,
while MENACE was one of the earliest physical implementations of a learning
agent for game playing. In this experiment we implement both ideas in Python
and observe how the behaviour of the agents changes over time.

\vspace{1em}

\noindent\textbf{Problem Description:}\\
The work is divided into two parts:
\begin{itemize}
  \item \textbf{Binary and 10-armed bandits:} The agent repeatedly chooses one
  of several actions (arms) and receives a stochastic reward. For the binary
  bandit, rewards are Bernoulli with fixed success probabilities. For the
  10-armed bandit, each arm's mean reward follows a random walk, making the
  problem non-stationary.
  \item \textbf{MENACE for Tic-Tac-Toe:} Game states correspond to board
  configurations, and each state has a ``matchbox'' containing beads for each
  legal move. Moves are sampled with probability proportional to the number of
  beads. After each game, bead counts are reinforced or penalised depending on
  whether MENACE wins, draws or loses.
\end{itemize}

\vspace{1em}

\noindent\textbf{Algorithm:}\\
For the bandit problems, we use the $\epsilon$-greedy strategy. With
probability $\epsilon$ the agent explores by choosing a random action; with
probability $1-\epsilon$ it exploits by choosing the action with the largest
estimated value $Q(a)$. In the stationary binary bandit the estimates are
updated using the sample average of observed rewards. For the non-stationary
10-armed bandit, a constant step-size $\alpha$ is used so that old rewards are
gradually forgotten.

For MENACE, each visited state--action pair is recorded during the game. At the
end of the game, bead counts are updated according to the outcome: winning
moves receive a larger positive increment, drawn games receive a small
increment, and losing moves are penalised but not allowed to go below one bead.
This simple rule gradually biases MENACE towards better strategies.

\vspace{1em}

\noindent\textbf{Implementation:}\\
The complete Python implementation (bandits and MENACE) for this experiment is
available at the following URL:\\[0.3em]
\hspace*{1em}\url{https://github.com/mkankotiya/AI-CS659/blob/main/search.py}

\vspace{1em}

\noindent\textbf{Results:}\\
In the stationary binary bandit, the $\epsilon$-greedy agent quickly learns to
prefer the arm with the higher success probability, and the average reward
converges close to the optimal value. In the non-stationary 10-armed bandit,
using a constant step-size allows the agent to track changing action values and
recover when the identity of the best arm changes.

For MENACE, the early games show many random and suboptimal moves. As more
games are played and bead updates accumulate, MENACE begins to avoid losing
lines and favour moves that lead to wins or draws. After sufficient training,
the engine rarely loses against a random opponent and often forces at least a
draw.

\vspace{1em}

\noindent\textbf{Conclusion:}\\
This experiment demonstrates how simple RL algorithms can learn effective
behaviour in both bandit settings and small games. The $\epsilon$-greedy
method illustrates the exploration--exploitation balance and the importance of
using appropriate updates for non-stationary environments. MENACE shows that
even a straightforward bead-based reinforcement scheme can produce competent
play in Tic-Tac-Toe. Together, these experiments provide an intuitive
introduction to reinforcement learning concepts that underlie more advanced
methods in artificial intelligence.

\end{document}
